% TOP_PROBLEM: {"problemId": "problem.four-exponentials-conjecture", "canonicalTitle": "Four Exponentials Conjecture", "releaseRank": 373, "primaryDomain": "number_theory_arithmetic_geometry", "primaryDomainLabel": "Number theory & arithmetic geometry", "displayStatus": "open", "releaseStatus": "open"}
% !TeX program = lualatex
% Definition and sources only; this file is not an LLM solution attempt.
\documentclass[11pt]{article}
\usepackage[margin=25mm]{geometry}
\usepackage{fontspec}
\setmainfont{Latin Modern Roman}
\usepackage{amsmath,amssymb,mathtools}
\usepackage{unicode-math}
\setmathfont{Latin Modern Math}
\usepackage{xurl,hyperref}
\hypersetup{colorlinks=true,urlcolor=blue}
\setcounter{secnumdepth}{0}
\setlength{\parindent}{0pt}
\setlength{\parskip}{5pt}
\setlength{\emergencystretch}{3em}
\begin{document}
\section*{\texorpdfstring{Four Exponentials Conjecture}{Four Exponentials Conjecture}}\label{373}
\subsection{Definitions and mathematical statement}
Let $x_1,x_2,y_1,y_2\in{\mathbb{C}}$ with $x_1,x_2$ linearly independent over ${\mathbb{Q}}$ and $y_1,y_2$ linearly independent over ${\mathbb{Q}}$. Thus each rational relation within either pair is trivial. The four exponentials conjecture asserts
\[
 \{e^{x_1y_1},e^{x_1y_2},e^{x_2y_1},e^{x_2y_2}\}
 \not\subset\overline{{\mathbb{Q}}}.
\]
Here $\overline{{\mathbb{Q}}}$ denotes the algebraic complex numbers, so at least one of the four values must be transcendental. The pairs need not consist of real numbers, and no independence between the two pairs is assumed. The assertion does not identify which value is transcendental or demand algebraic independence of the four values.
\par\smallskip\textit{Formulation and concept references:} \href{https://mathworld.wolfram.com/FourExponentialsConjecture.html}{[S1]}.

\subsection{Short English statement}
For two pairs of rationally independent complex numbers, must at least one exponential of a cross-product be transcendental?
\subsection{Sources}
\begin{itemize}
\item[S1]Four Exponentials Conjecture (MathWorld).
\url{https://mathworld.wolfram.com/FourExponentialsConjecture.html}.
\textit{Formulation source.}
\end{itemize}
\end{document}
